%!TEX root = ../../thesis.tex

\section{Overview}\label{sec:targets-overview}

% TODO: Open with the ladder and the table; the per-target sections then only add
% detail. The table is the most useful single page of this chapter --- one row per
% target, with columns:
%
%   target | capability isolated | formulation | reward source | role in evaluation
%
% Suggested contents (adjust to what is actually run):
%   high_and_low | single-step coverage choice        | input  | coverage    | sanity
%   <contextual> | state-dependent optimal action     | input  | coverage    | sanity
%   sequence     | deep credit assignment (16 deps)   | input  | designed    | diagnostic
%   branches     | planning past rewarding decoys     | input  | designed    | diagnostic
%   cjson        | grammar structure, raw coverage    | both   | raw         | headline
%   open62541    | protocol state, real stack         | input  | raw         | headline
%
% Then state the distinction the chapter is organized around, because it is the
% thesis's main defence against the "you designed the reward you then climbed"
% objection (\cref{sec:discussion-threats}):
%
%   - \cref{sec:targets-microbenchmarks}: rewards are hand-shaped and progress is
%     directly legible. These are instruments for diagnosis, not evidence about
%     fuzzing.
%   - \cref{sec:targets-realistic}: the reward is raw coverage and was not designed
%     around the agent. Real-world claims rest on these, and only on these.
%
% Say which formulation each target is evaluated under, and note that cjson is
% deliberately run under both --- that is what makes
% \cref{sec:results-head-to-head} a controlled comparison.
%
% Also list what every target shares, so it is not repeated six times: the coverage
% instrumentation of \cref{sec:platform-coverage}, the snapshot reset of
% \cref{sec:platform-snapshots}, and the frozen configuration of
% \cref{sec:methodology-hyperparameters}.
